\section{Lean definitions, proof targets, and verification boundaries}
\label{sec:lean}

\subsection{What has actually compiled}

The standalone project in \texttt{lean/} uses Lean 4.32.0, with no
Mathlib dependency. It contains 35 named theorem declarations and eight
checked examples across four Lean files. The successful build and
\texttt{\#print axioms} output are retained in
\texttt{lean/verification.log}. The files contain no admitted proof and
introduce no new axiom declaration.

\begin{center}
\begin{tabular}{p{0.24\textwidth}p{0.66\textwidth}}
\toprule
File & Checked content \\
\midrule
\texttt{Truncation.lean} & Rational clipping, clipping residual,
Lipschitz estimates, the bounded metric, the corrected weighted
truncation inequality, and the sharp excess-tail square bound. \\
\texttt{Energy.lean} & Binary and finite weighted energy identities,
telescoping, and an executable bounded search for a small energy
increment with its soundness and existence proof. \\
\texttt{Examples.lean} & A named counterexample to the unsafe
separate-tail shortcut, plus eight concrete exact examples. \\
\texttt{Audit.lean} & Axiom-dependency receipts for the principal
declarations. \\
\bottomrule
\end{tabular}
\end{center}

The central declaration
\texttt{truncation\_rho\_integrated} proves the repaired bound for an
arbitrary finite length and arbitrary nonnegative rational weights.
It is not merely a finite collection of examples. Similarly,
\texttt{energy\_window} produces a successful result of the executable
search \texttt{smallWindow}. Applied to the energy sequence along the
iterated challenge endpoints, it supplies the finite search used in
Theorem~\ref{m:metastability}.

These are host-Lean theorems. Their axiom audit lists
\texttt{propext}, \texttt{Classical.choice}, and
\texttt{Quot.sound}, inherited through the standard rational library
and proof infrastructure. They are not already proofs in
\(\Prf\,\mathrm{PRAAx}\), and we do not call them axiom-free or
choice-free formalizations. The human proofs use explicit rational
case splits and finite invariants, so their object-theory translation
has a clear route independent of these host axioms.

\subsection{The exact object-language endpoint}

The inspected canonical syntax is
\texttt{LeanProofs.ConciseIPRA.PTm}, with formulas
\texttt{Fm PTm n}. The intended judgment is
\[
 \Prf\,\mathrm{PRAAx}\,\Gamma\,\varphi,
 \qquad\text{written }\Gamma\vdash[\mathrm{PRAAx}]\varphi.
\]
The existing substitution, weakening, and cut interfaces are
\texttt{Prf.subst}, \texttt{Prf.weaken}, and \texttt{Prf.trans}.
The actual induction wrapper is \texttt{PRA.qfInd}; its step premise
has the induction formula added to the shifted context. The companion
\texttt{ObjectAdapter.lean.txt} records these exact inspected signatures
as an uncompiled integration candidate. It is excluded from the
standalone build.

We inspected existing object derivations for finite expectation,
signed parts, common denominators, finite kernels, Cauchy--Schwarz,
common-envelope truncation, block partition sums, error budgets, and
coherent presentations. Their precise source paths and hashes are
in \texttt{lean/source\_manifest.json}. We did not rebuild the full
repository, and do not turn old repository build receipts into a claim
of new verification.

\subsection{All data types needed for the next formalization stage}

The blueprint files specify the following data, with their coding and
proof fields kept separate.

\begin{description}
\item[Finite rational code.] A numerator pair \((p,n)\) and a
denominator \(d+1\), denoting \((p-n)/(d+1)\). All finite order
tests clear positive denominators.
\item[Dyadic step.] A level \(r\) and a vector of exactly \(2^r\)
finite rational codes. Product steps add a dimension and rectangular
level data. Refinement is explicit repetition and reindexing.
\item[Strong name.] A fixed PR producer of steps and the arithmetic
contract \(q\geq p\to\|a_q-a_p\|_1\leq2^{-p}\).
Program-code alternatives carry a fixed clock or supplied batch trace.
\item[Coherent mass presentation.] A fixed PR table
\(v(n,i)\), nonnegativity, and the equation
\(v(n,i)=v(n+1,2i)+v(n+1,2i+1)\). Derived programs produce the
density table and its exact rational energy.
\item[Quantitative tail data.] A fixed PR cutoff with a proved
first-order inequality. For Vitali it is \(U(p)\); for energy
convergence it is \(R(p)\). They are not arbitrary unproved function
fields whose existence is claimed by the framework.
\item[Bounded map.] A specified finite linear or bilinear map, a
norm bound, and its finite inequality. A pullback additionally records
partition compatibility and a compression bound.
\item[Formula with measured context.] Context dimension, atom
occurrence names in that context, rational connectives and their
sensitivities, integration, and permitted pullbacks. A shared
expression graph records its exact dependencies.
\item[Topological certificate.] Labelled rational distance table or
interval table, finite correspondence, parity-coded chains, scale,
and explicit prism output. Metric and chain equations are bounded
tests on the input data.
\item[Analytic input.] An \(L^1\) name, rational norm bounds, a
column-norm kernel name when required, and a positive denominator or
contraction gap. Complex rationals are pairs of rational codes.
\end{description}

The three accompanying blueprint files cover measure theory,
categorical and topological constructions, and analytic constructions.
Their pseudocode and Lean-shaped interfaces give the definitions and
theorem targets.
They are explicitly marked as blueprints. Their assumptions specify
missing proof obligations; they are not counted as established
theorems by being written as structure fields.

\subsection{Reifying the two central inequalities}

At a positive common denominator \(D\), let signed integers \(X,Y\)
represent the two values and let the natural threshold numerator be
\(K\), with \(D\leq K\). Define
\[
 A=|X-Y|,\quad R=\min(D,A),\quad
 T_X=\max(|X|-K,0),\quad T_Y=\max(|Y|-K,0).
\]
The denominator-cleared clipping statement is
\[
 D A\leq2KR+DT_X+DT_Y.
 \tag{C}\label{eq:clearedclip}
\]
The excess-square statement is
\[
                    4KT_X\leq X^2.
 \tag{S}\label{eq:clearedsquare}
\]
For \eqref{eq:clearedclip}, derive the three clipping-region identities
and nonexpansiveness by decidable signed-integer order cases. The
case split \(A\leq D\) or \(D\leq A\) bounds the clipped middle
term. The proof uses only semiring inequalities after moving negative
parts. For \eqref{eq:clearedsquare}, the nonzero branch has slack
\((|X|-2K)^2\). This is an explicit nonnegative polynomial
certificate.

To integrate, let \(w_i\) be nonnegative natural weights after common
denominator conversion. The induction invariant at prefix \(J\) is
\[
 D\sum_{i<J}w_i A_i
 \leq 2K\sum_{i<J}w_i R_i
       +D\sum_{i<J}w_i T_{X_i}+D\sum_{i<J}w_i T_{Y_i}.
\]
Replace each finite sum by its fixed PR symbol. The formula is
quantifier-free, its base is \(0\leq0\), and the successor step adds
one instance of \eqref{eq:clearedclip}. Existing expectation
monotonicity and linearity can implement this lifting. Modulus
composition then consists of term substitution and the proved rational
budget inequalities. No induction on ``all later terms converge'' is
needed.

For energy, derive the finite polynomial identity
\(2(a^2+b^2)=(a+b)^2+(a-b)^2\). Block-average coherence cancels the
cross term in the general identity. Existing regrouping and weighted
Cauchy--Schwarz supply the two bridges. The energy-window search
has the quantifier-free failed-search invariant
\(H_{t_J}\geq H_{t_0}+J\delta\); at the explicit bound this
contradicts the energy cap. This is the precise internal proof
architecture behind the new density and metastability constructions.

\subsection{Completion priorities}

The shortest useful next implementation is the object-theory version
of \eqref{eq:clearedclip}, its weighted lifting, and the rational
Vitali budget. That would close a concrete source gap and reuse the
existing arithmetic library. Next come the finite orthogonal-energy
identity and the local sibling bound. Completion and equality of names
then require only precision substitution and the specified rational
separation lemmas. The formula and Rips modules are independent
consumers of the same finite-array and error infrastructure. The
analytic applications come last, after the column-norm and moment
lemmas have actual object derivations.
